\chapter{Résultats, discussion et perspectives}
\label{ch:resultats}

\section{Fonctionnalités réalisées}

Le système livré comporte dix-huit modules fonctionnels, servis par 22 contrôleurs REST et
21 pages d'interface. Le tableau~\ref{tab:volumetrie} en donne la volumétrie, relevée sur
le commit final.

\begin{table}[H]
\centering
\caption{Volumétrie du système livré}
\label{tab:volumetrie}
\small
\begin{tabularx}{\textwidth}{@{}Xr@{\hskip 2em}Xr@{}}
\toprule
\textbf{Backend} & & \textbf{Frontend} & \\
\midrule
Entités JPA & 30 & Pages React & 21 \\
Services métier & 28 & Composants & 27 \\
Contrôleurs REST & 22 & Modules de design system & 8 \\
Repositories & 30 & Services d'appel API & 19 \\
Énumérations métier & 10 & & \\
Migrations Flyway & 3 & & \\
Tables relationnelles & 30 & & \\
Lignes de Java (sources) & 9\,980 & Lignes de JavaScript / JSX & 13\,471 \\
Lignes de Java (tests) & 1\,196 & Lignes SQL de migration & 750 \\
\bottomrule
\end{tabularx}
\end{table}

\section{Correspondance avec le cahier des charges}

Le tableau~\ref{tab:conformite} confronte chacune des 21 sections du cahier des charges à
son implémentation. La colonne « preuve » renvoie à la figure du rapport qui en atteste.

\begin{table}[H]
\centering
\caption{Correspondance entre le cahier des charges et l'implémentation}
\label{tab:conformite}
\scriptsize
\begin{tabularx}{\textwidth}{@{}c p{3.05cm} X p{1.85cm} c@{}}
\toprule
\textbf{§} & \textbf{Exigence} & \textbf{Implémentation} & \textbf{Preuve} & \textbf{État} \\
\midrule
1 & Objectifs du système & Architecture trois tiers, 18 modules & Fig.~\ref{fig:archi} & Couvert \\
2.1 & Accès visiteur libre & Endpoints publics, portail & Fig.~\ref{fig:portail} & Couvert \\
2.2 & Cinq statuts de membre & \code{MemberStatus}, badges dédiés & Fig.~\ref{fig:membres} & Couvert \\
2.2.C & Doctorant sans modules internes & Exclusion au niveau classe & Fig.~\ref{fig:sidebarPhd} & Couvert \\
3 & Profil complet et modifiable & \code{Member}, 11 champs, avatar & Fig.~\ref{fig:membreGouv} & Couvert \\
4.1 & Date de création du laboratoire & \code{Laboratory} & Fig.~\ref{fig:admin} & Couvert \\
4.2 & Dates du membre & 3 champs datés, naissance confidentielle & Fig.~\ref{fig:membreHist} & Couvert \\
4.3 & Historique des affiliations & \code{AffiliationHistory}, retour = nouvelle période & Fig.~\ref{fig:membreHist} & Couvert \\
5.1 & Administrateur système & Contrôleur réservé, identité du labo & Fig.~\ref{fig:admin} & Couvert \\
5.2 & Directeur en mandat & Décisions réservées au directeur & Fig.~\ref{fig:demandes} & Couvert \\
5.3 & Mandats datés, historique conservé & \code{Mandate}, clôture sans suppression & Fig.~\ref{fig:gouvMandats} & Couvert \\
6 & Demandes d'adhésion & Dépôt public, décision du directeur, création de compte & Fig.~\ref{fig:apply}, \ref{fig:demandes} & Couvert \\
7 & Rôles et responsabilités & \code{UserRoleSyncService}, rôle dérivé et historisé & Fig.~\ref{fig:seqMandat} & Couvert \\
8.1 & Retraité : compte gelé & Gel vérifié à chaque requête & Fig.~\ref{fig:etats} & Couvert \\
8.2 & Ancien membre réactivable & \code{returnToLab()} & Fig.~\ref{fig:etats} & Couvert \\
9 & Événements, archivage par édition & \code{editionNumber}, organisateurs, participants & Fig.~\ref{fig:events} & Couvert \\
10 & Archivage documentaire versionné & \code{DocumentVersion}, 9 catégories & Fig.~\ref{fig:docVersions} & Couvert \\
11 & Rapport annuel automatique & Agrégation, PDF, planification annuelle & Fig.~\ref{fig:rapports} & Couvert \\
12 & Publications, doctorants inclus & Classement auteur / année / équipe & Fig.~\ref{fig:publis} & Couvert \\
13 & Communication interne & 3 portées de discussion & Fig.~\ref{fig:notifs} & Couvert \\
14 & Matériel : reçu / non reçu & Requête de complément dédiée & Fig.~\ref{fig:equipDist} & Couvert \\
15 & Besoins en matériel & Soumission, validation par la direction & Fig.~\ref{fig:equip} & Couvert \\
16 & Conventions et collaborations & \code{Partner}, \code{Agreement} & Fig.~\ref{fig:conventions} & Couvert \\
17 & Procès-verbaux & Réunions, ordre du jour, dépôt de PV & Fig.~\ref{fig:reunions} & Couvert \\
18 & Recherche globale & Endpoint transverse, barre permanente & Fig.~\ref{fig:dashboard} & Couvert \\
19 & Notifications & 4 politiques de destinataires & Fig.~\ref{fig:notifs} & Couvert \\
20 & Calendrier & Agrégation événements + réunions & Fig.~\ref{fig:calendrier} & Couvert \\
21 & Règles générales & Suppression logique, 4 historiques, audit & Fig.~\ref{fig:audit} & Couvert \\
\bottomrule
\end{tabularx}
\end{table}

Les 21 sections sont couvertes par une fonctionnalité opérationnelle. Cette conformité
porte sur le périmètre fonctionnel ; elle ne signifie pas que le système est exempt de
limites, exposées en~§\ref{sec:limites}.

\section{Points forts}

\paragraph{La séparation \srt{}~/~\rol{}~/~\man{} est structurelle, pas déclarative.} Elle
n'est pas seulement documentée : elle est imposée par le code. Le rôle applicatif ne peut
pas être incohérent avec la réalité institutionnelle, parce qu'il en est calculé. La
réconciliation au démarrage garantit en outre que cette propriété reste vraie même si des
données arrivent par un autre chemin.

\paragraph{La sécurité ne fait pas confiance au jeton.} Le rechargement systématique des
habilitations depuis la base rend la révocation immédiate. C'est un choix coûteux en
lectures mais qui supprime toute une classe de vulnérabilités liées aux jetons survivants.

\paragraph{L'interface et le serveur appliquent la même matrice.} Les listes de rôles
utilisées par les gardes de route et par le masquage des commandes correspondent aux
annotations du serveur. L'audit a permis de corriger le seul écart constaté ; l'objectif
--- aucune commande visible qui soit refusée, aucune opération autorisée qui soit masquée
--- est atteint.

\paragraph{Rien n'est jamais détruit.} La suppression logique généralisée, les quatre tables
d'historique et le journal d'audit répondent ensemble à la règle §21 du cahier des charges.
L'utilitaire de résolution sûre des associations archivées évite que cette conservation ne
dégrade l'expérience.

\paragraph{La validation ne se limite pas à la compilation.} Les défauts décrits en
§\ref{sec:defauts} n'auraient été révélés ni par la compilation, ni par une lecture de
code, ni par la vérification des codes de statut HTTP. Ils l'ont été par un parcours réel
et par des tests d'intégration.

\section{Limites}
\label{sec:limites}

Les limites suivantes sont assumées et exposées sans atténuation.

\paragraph{Absence de thème sombre.} Le document de spécification le mentionne parmi les
exigences non fonctionnelles. Il n'a pas été implémenté. Le système de composants repose sur
des variables de couleur centralisées, ce qui rendrait l'ajout mécanique, mais le travail
n'a pas été fait.

\paragraph{Aucune mesure de couverture de code.} Aucun outil de couverture n'est configuré.
Les 60 tests protègent les règles identifiées comme critiques, mais la proportion du code
exercée n'est pas connue et n'est pas estimée dans ce rapport.

\paragraph{Modules inégalement couverts par les tests.} Les discussions, les partenariats et
la génération PDF ne disposent pas de tests automatisés dédiés. Ils ont été validés
manuellement lors de l'audit, ce qui est un niveau de garantie inférieur.

\paragraph{Pas de campagne de performance.} Les précautions usuelles sont prises ---
pagination côté serveur, désactivation de la session étendue à la vue, index sur les
colonnes de recherche --- mais aucune mesure sous charge n'a été conduite. Le comportement
du système au-delà de quelques dizaines d'utilisateurs simultanés n'est pas documenté.

\paragraph{Accessibilité partielle.} Les boutons-icônes portent des libellés accessibles, le
focus est visible et les modales se ferment au clavier, mais aucun audit selon un référentiel
formel n'a été mené.

\paragraph{Rafraîchissement de session non automatisé.} L'expiration du jeton d'accès ramène
l'utilisateur à l'écran de connexion ; le jeton de rafraîchissement existe et l'endpoint
correspondant est exposé, mais le client ne l'utilise pas de manière transparente.

\paragraph{Client livré en un seul lot.} Le paquet JavaScript n'est pas découpé en fragments
chargés à la demande. Sur une connexion lente, le premier affichage en est pénalisé.

\paragraph{Valeur de repli de la clé de signature.} Le fichier de configuration conserve
une clé de démonstration utilisée à défaut de variable d'environnement. Elle est surchargée
partout dans le projet et porte la mention \code{change-in-production}, mais un
déploiement distrait hériterait d'une clé publique. La rendre obligatoire est une
modification de deux lignes, repérée et non appliquée : elle touche la configuration des
tests d'intégration, et le périmètre de la dernière passe de correction était limité au
compte d'amorçage.

\paragraph{Historique de gestion de version à assainir.} Comme exposé au
chapitre~\ref{ch:securite}, un identifiant de base de données a figuré dans un fichier
versionné. Il a été retiré du code et le mot de passe correspondant a été effectivement
changé --- l'ancien est refusé par le serveur de base de données. L'historique du dépôt en
conserve néanmoins la trace, désormais inopérante ; sa purge reste à effectuer.

\section{Perspectives d'évolution}

\paragraph{Court terme.} Ajouter le thème sombre, configurer un outil de couverture pour
objectiver l'effort de test, étendre les tests automatisés aux trois modules qui en sont
dépourvus, et mettre en œuvre le rafraîchissement transparent du jeton.

\paragraph{Moyen terme.} Découper le paquet client par route, introduire une recherche
plein texte indexée sur les documents et les publications, et exporter le rapport annuel
dans un second format destiné à la réutilisation bureautique.

\paragraph{Long terme, en lien avec la spécialité.} Les données accumulées par la plateforme
--- historiques de statuts, de mandats, d'affiliations, production scientifique datée,
journal d'audit --- constituent un jeu longitudinal exploitable. Une extension analytique
naturelle consisterait à en dériver des indicateurs de pilotage : évolution des effectifs
par statut, productivité des équipes, durée moyenne des mandats, saisonnalité de
l'activité. La conception actuelle s'y prête, précisément parce qu'aucune donnée n'est
détruite : chaque table d'historique est déjà une série temporelle. C'est l'articulation la
plus directe entre ce projet de module JEE et la finalité du Master.
